\documentclass[11pt]{article}

\usepackage[T1]{fontenc}
\usepackage{lmodern}
\usepackage{amsmath,amssymb}
\usepackage{fullpage}
\usepackage{setspace}

\setlength{\parindent}{0pt}

\newcommand{\T}{\mathbb{T}}

\begin{document}

\onehalfspacing

This is very surprising. At first, I assumed that there must be a mistake
somewhere: this resolves the strongest form of the question with an argument
only a page long, using completely standard harmonic analysis. The second,
weaker question was resolved by Beck in a 44-page \emph{Annals} paper with a
much more analytically involved proof.

I now believe, however, that the proof is correct, though I encourage others
to read it with a sceptical eye as well. This seems to be a question for which
attacking the more difficult version somehow makes the problem easier---but
going from a 44-page \emph{Annals} paper to a single page of elementary
harmonic analysis is a dramatic jump!

The proof is short enough to give in its entirety below.

For convenience, fix $n$ and the finite truncation $z_1,\ldots,z_n$. Write
\[
  \T:=\mathbb{R}/\mathbb{Z},
  \qquad
  e(x):=e^{2\pi i x},
  \qquad
  z_j=e(x_j),
\]
and let
\[
  f_k:=\sum_{j\leq k}\delta_{x_j},
  \qquad
  f:=f_n.
\]
Define
\[
  M_k:=\max_{x\in\T}\prod_{j\leq k}|e(x)-e(x_j)|.
\]

As usual, we first take logarithms. Let
\[
  B(x):=\log|2\sin(\pi x)|,
  \qquad
  F_k:=B*f_k,
  \qquad
  L_k:=\sup_{x\in\T}F_k(x)=\log M_k.
\]
With the usual Fourier-series normalization,
\[
  \widehat{B}(0)=0,
  \qquad
  \widehat{B}(m)=-\frac{1}{2|m|}
  \quad (m\neq 0).
\]

Let $H$ be chosen later, and let $K_H$ be the Fej\'er kernel, normalized so
that $K_H\geq 0$ and $\int_{\T}K_H=1$. For every $y\in\T$,
\[
  (K_H*F_k)(y)\leq L_k.
\]
Taking $y=x_{k+1}$ and summing over $1\leq k<n$ gives
\[
  \sum_{k<n}L_k
  \geq
  \sum_{1\leq i<j\leq n}(B*K_H)(x_j-x_i).
\]
Since $L_n\geq 0$, the same lower bound holds with
$\sum_{k\leq n}L_k$ on the left.

Expanding the right-hand side in a Fourier series gives
\[
  \sum_{1\leq i<j\leq n}(B*K_H)(x_j-x_i)
  =
  \frac{n}{2}\sum_{m\leq H}\frac{\widehat{K_H}(m)}{m}
  -
  \frac{1}{2}\sum_{m\leq H}
  \frac{\widehat{K_H}(m)}{m}|\widehat{f}(m)|^2.
\]
Now
\[
  \widehat{K_H}(m)=1+O\left(\frac{m}{H}\right)
  \qquad (1\leq m\leq H),
\]
and therefore
\[
  \sum_{k\leq n}\log M_k
  \geq
  \frac{n}{2}\log H
  -
  O\left(
    n+\sum_{m\leq H}\frac{|\widehat{f}(m)|^2}{m}
  \right).
\]

We use the following elementary observation. If $F$ is real-valued,
$\int_{\T}F=0$, and $C:=\sup_{\T}F$, then
\[
  |\widehat{F}(m)|\leq C
  \qquad (m\neq 0).
\]
Indeed, $G:=C-F\geq 0$, so for $m\neq 0$,
\[
  |\widehat{F}(m)|
  =
  |\widehat{G}(m)|
  \leq
  \|G\|_1
  =
  \int_{\T}G
  =
  C.
\]

Apply this observation to $F_n=B*f$. If $L_n>2\log n$, then
$M_n>n^2$, which is already stronger than the bound we seek. We may
therefore assume that $L_n\leq 2\log n$. Since
\[
  \widehat{F_n}(m)
  =
  -\frac{\widehat{f}(m)}{2|m|}
  \qquad (m\neq 0),
\]
we obtain
\[
  |\widehat{f}(m)|\ll m\log n.
\]
Consequently,
\[
  \sum_{k\leq n}\log M_k
  \geq
  \frac{n}{2}\log H
  -
  O\bigl(n+(H\log n)^2\bigr).
\]

Choosing
\[
  H\asymp \frac{\sqrt{n}}{\log n}
\]
and applying AM--GM, we conclude that
\[
  \frac{1}{n}\sum_{k\leq n}M_k
  \geq
  \exp\left(\frac{1}{n}\sum_{k\leq n}\log M_k\right)
  \gg
  \frac{n^{1/4}}{\sqrt{\log n}}.
\]

The new ``magic'' step, I think, is the bound
\[
  L_k\geq (K_H*F_k)(x_{k+1}).
\]
The same inequality holds with any evaluation point, but the choice
$x_{k+1}$ means that averaging over $k$ produces the particularly
well-behaved pair energy
\[
  \sum_{1\leq i<j\leq n}(B*K_H)(x_j-x_i).
\]
Thus we pass from a sum of upper extrema, which is awkward to handle, to a
quadratic energy that is transparent on the Fourier side. This trick seems so
natural and powerful that surely it has appeared in other arguments in
harmonic analysis and must have other applications.

\end{document}
